\documentclass[11pt,a4paper]{article}
\usepackage[T1]{fontenc}
\usepackage{lmodern,amsmath,amssymb,amsthm,mathtools,microtype}
\ifdefined\pdfmapfile\pdfmapfile{+lm.map}\pdfmapfile{+cm.map}\pdfmapfile{+symbols.map}\fi
\usepackage[margin=25mm,headheight=14pt]{geometry}
\usepackage{xcolor,booktabs,array,longtable,listings,fancyhdr}
\definecolor{ink}{RGB}{24,54,75}
\usepackage[colorlinks=true,linkcolor=ink,citecolor=ink,urlcolor=ink]{hyperref}
\hypersetup{pdftitle={Three Proof Transfers: Graph Spectra, Bessel Ratios, and Multiplicative Density},pdfauthor={Research prepared for Keio}}
\pagestyle{fancy}\fancyhf{}
\fancyhead[L]{\small Three proof transfers}
\fancyhead[R]{\small Unified research manuscript}
\fancyfoot[C]{\thepage}
\newtheorem{theorem}{Theorem}[section]
\newtheorem{lemma}[theorem]{Lemma}
\newtheorem{corollary}[theorem]{Corollary}
\newtheorem{proposition}[theorem]{Proposition}
\theoremstyle{remark}\newtheorem{remark}[theorem]{Remark}
\newcommand{\R}{\mathbb R}
\newcommand{\N}{\mathbb N}
\newcommand{\Z}{\mathbb Z}
\newcommand{\C}{\mathcal C}
\newcommand{\dens}{\operatorname{d}}
\newcommand{\tr}{\operatorname{tr}}
\newcommand{\one}{\mathbf 1}
\numberwithin{equation}{section}
\setlength{\emergencystretch}{2em}
\allowdisplaybreaks[1]
\raggedbottom
\setcounter{tocdepth}{1}
\lstset{language=Python,basicstyle=\ttfamily\footnotesize,
  breaklines=true,columns=fullflexible,keepspaces=true,
  showstringspaces=false,frame=single,rulecolor=\color{black!20},
  xleftmargin=4pt,xrightmargin=4pt,aboveskip=9pt,belowskip=9pt}
\title{\textbf{Three Proof Transfers}\\[6pt]
\large Graph Spectra, Bessel Ratios, and Multiplicative Density}
\author{Research prepared for Keio}
\date{15 November 2026\\\small Unified analytic proof draft for independent review}
\begin{document}
\maketitle
\begin{abstract}
We assemble three complete arguments into one manuscript with common
notation conventions, section-based theorem numbering, and a single bibliography.
For connected simple graphs, the square-energy theorem of Liu, Tang and
Zhang yields the strict Randi\'c bound $R>\lambda_2+1$, its two edge-count
strengthenings, and a quadratic estimate with classified equality cases.
For modified Bessel functions, a weighted Wronskian gives a complete
classification of monotonicity of $I_{\nu_1}(a_1x)/I_{\nu_2}(a_2x)$
for $\nu_1\ge\nu_2>-1$ and $0<a_1\le a_2$; a totally positive
integral transform transfers unimodality to Nuttall $Q$ ratios.
The distinct log-concavity clause is false and is treated separately.
For every pairwise coprime integer triple $2\le a<b<c$, a nonnegative
packing of directed lattice corners determines the largest natural
density of a set avoiding $\{x,ax,bx,cx\}$ as
$1-\min\{(ab+c)/((a+1)(b+1)(c+1)),1/(c+1)\}$.
All original mathematical proofs are retained. External theorems are
explicitly attributed; the two finite-check programs are included as
supplements, not as replacements for proof. This manuscript does not
assert historical priority, external peer review, or proof-assistant verification.
\end{abstract}
\tableofcontents

\section{Introduction and common conventions}
Each argument transfers a tractable comparison into a global conclusion.
The graph argument combines squared spectral energy with a degree-weighted
Rayleigh bound. The special-function argument converts a one-change sign
pattern for a Wronskian derivative into unimodality and then applies a
totally positive transform. The arithmetic argument turns an explicit
weighted packing into a sharp logarithmic estimate and an attained
natural density. These are three independent results; no claim that one
of their main theorems implies either of the others is needed.

Section~\ref{sec:graph} addresses Cazenave's 2026 graph conjecture using
the explicitly cited square-energy theorem. Section~\ref{sec:bessel}
addresses the unimodality clause of the 2025 Karp--Vishnyakova--Zhang
conjecture and retains all parameter boundary cases.
Section~\ref{sec:density} proves an exact density formula, without
identifying it as the resolution of a previously named conjecture.
Thus the original scope and provenance of the three source manuscripts
remain distinct in this unified edition.

Notation is local to each main section when a symbol has different
standard uses: $R$ denotes the Randi\'c index in Section~\ref{sec:graph},
$R(x)$ a Bessel ratio in Section~\ref{sec:bessel}, and $R(z)$ a packing
tail in Section~\ref{sec:density}. All theorem-like environments share
one counter within each section; equations and cross-references are
numbered consistently throughout.

The source manuscripts are \emph{A Square-Energy Route to Cazenave's
Randi\'c Conjecture}, \emph{Unimodality of Modified Bessel Ratios and
Nuttall $Q$ Ratios}, and \emph{Exact Natural Densities for Multiplicative
Four-Point Configurations}, all dated 15 November 2026. Their full
proofs and substantive qualifications are preserved below. Repeated
front matter and bibliographies have been consolidated. The source
files used for this edition were recovered from their original three-paper bundle.

\section{Graph spectra and the Randi\'c index}
\label{sec:graph}
\subsection{Statements and scope}
Let $G$ be a finite connected undirected simple graph, with $n\ge2$
vertices, $m$ edges, degrees $d_v$, adjacency matrix $A$, and adjacency
eigenvalues $L=\lambda_1\ge t=\lambda_2\ge\cdots\ge\lambda_n$.
All edges below are unordered and are counted once. Set
\[
 R=\sum_{uv\in E(G)}\frac1{\sqrt{d_ud_v}},\qquad
 s^\pm=\sum_{\pm\lambda_i>0}\lambda_i^2.
\]
Cazenave \cite[Section~7, Conjecture~1, \#643]{graph-Caz} states
\begin{quote}$R\ge\lambda_2+1$.\end{quote}
The same source establishes the bipartite case and leaves the
non-bipartite case open. It also proposes the stronger assertions
$2m\ge L^2+(t+1)^2$ and $m\ge L(t+1)$.

\begin{theorem}\label{graph-thm:main}
Every connected simple graph on at least two vertices satisfies
\[
 2m>L^2+(t+1)^2,\qquad m>L(t+1),\qquad R>t+1.
\]
If the graph is not complete, then the quantitative estimate
\begin{equation}\label{graph-eq:gap}
 R-t-1\ \ge\
 \frac{(L-t-1)^2+n-2t-2}{2L}\ >0
\end{equation}
also holds.
\end{theorem}

\begin{theorem}\label{graph-thm:square}
Every connected noncomplete graph satisfies
\begin{equation}\label{graph-eq:square}
                R^2\ge t^2+n-1.
\end{equation}
Equality holds precisely for stars $K_{1,n-1}$, with $n\ge3$.
Consequently, for every connected graph with $n\ge4$, inequality
\eqref{graph-eq:square} holds, with equality precisely for stars and $K_4$.
The excluded complete graphs $K_2,K_3$ do not satisfy \eqref{graph-eq:square}.
\end{theorem}

The new deduction below uses the 2026 theorem of Liu--Tang--Zhang,
not a newly established square-energy theorem. Its proof and attribution
are kept separate. VMRC manuscript~[161] \cite{VMRC} suggested comparing
the same degree-weighted invariant to a global graph parameter; its
radius theorem is not a premise here. The graph result concerns an adjacency eigenvalue; the cited corpus
connection is methodological, not a theorem assumed in the proof.

\subsection{Three ingredients}
\begin{lemma}[Strict sign-support bound]\label{graph-lem:nodal}
For a connected graph with $t\ge0$ and $n\ge3$,
\[
                          t<\frac n2-1.
\]
\end{lemma}
\begin{proof}
The Perron vector is strictly positive. An eigenvector $y$ for $t$,
orthogonal to that vector, therefore has both positive and negative
coordinates. Write $P=\{v:y_v>0\}$, $N=\{v:y_v<0\}$, and
$p=|P|$, $q=|N|$.
At a vertex $u$ maximizing $y$, the eigenvalue equation gives
\[
 t y_u=\sum_{v\sim u}y_v
       \le\sum_{v\in P\cap N(u)}y_v\le(p-1)y_u.
\]
Applying the same argument to $-y$ gives $t\le q-1$.
Thus $n\ge p+q\ge2(t+1)$.

If equality $n=2(t+1)$ held, then $p=q=t+1$, there would be no
zero coordinates, and both inequalities at $u$ would be equalities.
Consequently $u$ would be adjacent to every other vertex of $P$,
and all coordinates on $P$ would equal $y_u>0$.
At any vertex of $P$, attaining the value $(p-1)y_u$ in its
eigenvalue equation then requires all $p-1$ possible neighbours in
$P$ and no neighbours in $N$. Thus no edge joins $P$ and $N$,
contrary to connectedness. This also covers $p=1$: the equation
with $t=0$ cannot allow a negative neighbour.
\end{proof}

\begin{lemma}[Favaron--Mah\'eo--Sacl\'e]\label{graph-lem:fav}
One has $m\le LR$. If equality holds, $d_ud_v$ is constant over all
edges. A connected graph with this constant-product property is either
regular or bipartite with a constant degree on each part.
\end{lemma}
\begin{proof}
Apply Rayleigh's principle to $x_v=\sqrt{d_v}$:
\[
 S:=\sum_{uv\in E}\sqrt{d_ud_v}\le mL.
\]
Cauchy--Schwarz gives $m^2\le SR\le mLR$, proving the bound.
Equality in Cauchy--Schwarz forces $(d_ud_v)^{1/2}$ to be constant
on edges. Along a path the degrees then alternate between two values.
If there is an odd cycle those values agree and connectedness makes
the graph regular. Otherwise the graph is bipartite and the degrees
are constant separately on its two parts.
\end{proof}

\begin{lemma}[Liu--Tang--Zhang, Theorem~1.2]\label{graph-lem:energy}
For every connected graph,
\[
                  \min(s^+,s^-)\ge n-1.
\]
\end{lemma}
This is the principal external theorem \cite{graph-LTZ}.
The original source supplies a matrix proof of this theorem. Its
hypotheses are exactly
finite, simple, undirected and connected; no bipartiteness,
regularity or degree restriction is imposed.

\subsection{Proof of the strict conjecture}
\begin{proof}[Proof of Theorem~\ref{graph-thm:main}]
If $G$ is complete, then $m=n(n-1)/2$, $L=n-1$, $t=-1$ and $R=n/2$.
All three strict inequalities follow directly.

Suppose $G$ is not complete. A shortest path between two nonadjacent
vertices contains an induced $P_3$. Eigenvalue interlacing therefore
gives $t\ge\lambda_2(P_3)=0$; in particular $n\ge3$.
Lemma~\ref{graph-lem:energy} and $s^++s^-=2m$ give
\begin{equation}\label{graph-eq:twoenergy}
                2m\ge L^2+t^2+n-1.
\end{equation}
By Lemma~\ref{graph-lem:nodal}, $n-1>2t+1$. Hence
\[
 2m>L^2+(t+1)^2\ge2L(t+1).
\]
Since $L>0$ and $LR\ge m$, the other two assertions follow.
For the quantitative form, subtract $t+1$ from
$R\ge m/L\ge(L^2+t^2+n-1)/(2L)$ and complete the square.
The final numerator in \eqref{graph-eq:gap} is strictly positive because
$n-2t-2>0$.
\end{proof}

When $n$ is odd and the graph is noncomplete, the sign-support proof
also gives $t\le\min(p,q)-1\le\lfloor n/2\rfloor-1$. Therefore
\eqref{graph-eq:gap} yields the explicit order-parity margin
\[
 R-t-1\ge\frac{(L-t-1)^2+1}{2L}\ge\frac1{2L}.
\]

\subsection{The quadratic strengthening and its equality cases}
\begin{proof}[Proof of Theorem~\ref{graph-thm:square}]
For a noncomplete graph, \eqref{graph-eq:twoenergy} implies
\[
 R\ge\frac mL\ge\frac{L^2+t^2+n-1}{2L}
                \ge\sqrt{t^2+n-1},
\]
where the last step is the arithmetic--geometric mean inequality.

Suppose equality holds. Equality throughout this chain forces
\begin{equation}\label{graph-eq:equality}
 R=L,\qquad m=L^2,\qquad L^2=t^2+n-1,
 \qquad m=LR.
\end{equation}
By Lemma~\ref{graph-lem:fav}, there are two possibilities.

If the graph is regular of degree $d$, then $L=d$ and $m=nd/2$.
Thus $n=2d$ and $t^2=(d-1)^2$. Here $d\ge2$ and $t\ge0$, so
$t=d-1=n/2-1$, contradicting Lemma~\ref{graph-lem:nodal}.

Otherwise the graph is bipartite, with degrees $a$ on one part and
$b$ on the other. Its positive eigenvector taking values $\sqrt a$
and $\sqrt b$ on the respective parts shows $L=\sqrt{ab}$.
Since $m=ab$, the part sizes are $b$ and $a$; every possible edge
between the parts is present. Thus $G=K_{a,b}$. As $n\ge3$, its
second eigenvalue is $t=0$, so \eqref{graph-eq:equality} gives
$ab=a+b-1$, or $(a-1)(b-1)=0$. The graph is a star.
Conversely, a star has $t=0$ and $R^2=n-1$.

For $K_n$, \eqref{graph-eq:square} becomes $n^2/4\ge n$, which holds
exactly when $n\ge4$ and is an equality exactly when $n=4$.
This proves the final assertions.
\end{proof}

There is also a version involving every nonprincipal positive
eigenvalue. The same argument, replacing $t^2$ by
$P:=s^+-L^2\ge0$, proves for every connected graph that
\[
                         R^2\ge n-1+P.
\]
The two-eigenvalue form was singled out because it resolves the
stated conjecture and has the elementary equality classification above.

\subsection{Sharpness and research status}
Let $B_k$ consist of two copies of $K_k$ joined by one edge, $k\ge2$.
The degree count and the symmetric/antisymmetric quotient matrices give
\[
 R(B_k)=k-2+2\sqrt{1-\frac1k}+\frac1k,
 \qquad
 \lambda_2(B_k)=\frac{k-3+\sqrt{(k+1)^2-4}}2.
\]
For completeness, the two quotient matrices are
\[
 \begin{pmatrix}1&k-1\\1&k-2\end{pmatrix},\qquad
 \begin{pmatrix}-1&k-1\\1&k-2\end{pmatrix}.
\]
Their positive roots are the largest and second-largest adjacency
eigenvalues; all remaining eigenvalues are negative (including the
eigenvalue $-1$ on within-clique zero-sum subspaces).
Consequently $R(B_k)-\lambda_2(B_k)-1\to0$ as $k\to\infty$.
Thus the strict inequality cannot be improved by any fixed positive
additive constant uniformly over connected graphs.

The literature check identified Conjecture~1 in Cazenave's original
preprint and the square-energy theorem used above. Searches also found
solutions of other numbered Cazenave conjectures, which are different
statements. The check did not establish whether the particular deduction
in this section has already appeared elsewhere. No historical-priority,
peer-review or machine-formalization claim is made for this section.
The proof itself is analytic and finite-dimensional; enumerating small
graphs is not an input.

\section{Modified Bessel ratios and Nuttall \texorpdfstring{$Q$}{Q} ratios}
\label{sec:bessel}
\subsection{The original question and the precise result}
For real $\nu>-1$ and $z>0$, the modified Bessel function is
\begin{equation}\label{bessel-eq:series}
 I_\nu(z)=\sum_{k=0}^{\infty}
 \frac{(z/2)^{2k+\nu}}{k!\Gamma(k+\nu+1)}>0.
\end{equation}
Karp--Vishnyakova--Zhang formulate two assertions in
\cite[Conjecture~1]{bessel-KVZ}: unimodality of
$I_{\nu_1}(a_1x)/I_{\nu_2}(a_2x)$ for
$\nu_1\ge\nu_2>-1$, $0<a_1\le a_2$, and log-concavity when
$\nu_2>0$. The parameter $b$ printed in that conjecture is absent
from this Bessel ratio. It becomes relevant to their preceding
Nuttall $Q$ application.

Here unimodality means that the function is nondecreasing and then
nonincreasing, allowing either part to be empty. This implies the
authors' convention of at most one change of monotonicity.

\begin{theorem}\label{bessel-thm:main}
Suppose $\nu_1\ge\nu_2>-1$ and $0<a_1\le a_2$. Put
\[
 R(x)=\frac{I_{\nu_1}(a_1x)}{I_{\nu_2}(a_2x)},\qquad x>0.
\]
The following cases exhaust all possibilities:
\begin{enumerate}
\item If $\nu_1=\nu_2$ and $a_1=a_2$, then $R\equiv1$.
\item If $\nu_1=\nu_2$ and $a_1<a_2$, then $R$ is strictly decreasing.
\item If $\nu_1>\nu_2$ and $a_1<a_2$, then there is exactly one
$x_*>0$ such that $R$ strictly increases on $(0,x_*)$ and strictly
decreases on $(x_*,\infty)$.
\item If $\nu_1>\nu_2$, $a_1=a_2$, and $\nu_1+\nu_2\ge0$,
then $R$ is strictly increasing.
\item If $\nu_1>\nu_2$, $a_1=a_2$, and $\nu_1+\nu_2<0$,
then $R$ has exactly one strict maximum as in case~3.
\end{enumerate}
In particular, the unimodality assertion in \cite[Conjecture~1]{bessel-KVZ}
holds throughout its stated parameter range.
\end{theorem}

The theorem proves the unimodality clause, not the false conjunction
of unimodality and unrestricted positive-order log-concavity.
The literature check verified the original 2025 source and examined
later indexed discussions and related Bessel-ratio literature, including
\cite{bessel-Segura}. It did not establish worldwide priority. An explicit
formal counterexample to the log-concavity clause is already recorded
in \cite{bessel-TT}; the negative conclusion below is not claimed as new.
The proof relies on classical Bessel identities, not on the correctness
of any unpublished corpus theorem.

\subsection{A weighted Wronskian identity}
Write
\[
 f(x)=I_{\nu_1}(a_1x),\quad g(x)=I_{\nu_2}(a_2x),\quad
 d=\nu_1-\nu_2,\quad s=\nu_1+\nu_2,
\]
and put $B=ds=\nu_1^2-\nu_2^2$, $D=a_2^2-a_1^2\ge0$.
Both $f$ and $g$ are strictly positive. The modified Bessel equation
\cite[Eq.~10.25.1]{bessel-DLMF} gives
\[
 f''+\frac1x f'=\left(a_1^2+\frac{\nu_1^2}{x^2}\right)f,
 \qquad
 g''+\frac1x g'=\left(a_2^2+\frac{\nu_2^2}{x^2}\right)g.
\]
Define the weighted Wronskian, with the indicated sign convention,
\begin{equation}\label{bessel-eq:W}
 W(x)=x\bigl(f'(x)g(x)-f(x)g'(x)\bigr).
\end{equation}
Direct differentiation, followed by cancellation of the first-derivative
terms, yields the exact identities
\begin{equation}\label{bessel-eq:identity}
 \boxed{\quad W'(x)=\left(-Dx+\frac{B}{x}\right)f(x)g(x),
 \qquad R'(x)=\frac{W(x)}{xg(x)^2}.\quad}
\end{equation}
Thus the sign of $R'$ is the sign of $W$, while the sign of $W'$ is
controlled by the elementary function $B-Dx^2$.

\begin{lemma}[Endpoint information]\label{bessel-lem:endpoints}
If $d>0$, there is $C>0$ such that, as $x\downarrow0$,
\begin{equation}\label{bessel-eq:smallW}
 W(x)=Cd\,x^s\bigl(1+O(x^2)\bigr)>0
 \quad\hbox{for all sufficiently small }x>0.
\end{equation}
If $d=0$, then $W(x)\to0$ as $x\downarrow0$.
When $a_1<a_2$, $W(x)<0$ for all sufficiently large $x$.
When $a_1=a_2=a$ and $B\ne0$, the eventual sign of $W$ is the sign
of $B$.
\end{lemma}
\begin{proof}
The series \eqref{bessel-eq:series}, with its termwise derivative, gives
\[
 I_\nu(ax)=\frac{(a/2)^\nu}{\Gamma(\nu+1)}x^\nu
 \left(1+\frac{a^2x^2}{4(\nu+1)}+O(x^4)\right).
\]
Consequently
\[
 \frac{f'}f-\frac{g'}g=\frac d x+O(x),\qquad
 fg=Cx^s\bigl(1+O(x^2)\bigr),
\]
which proves \eqref{bessel-eq:smallW}. If $d=0$, the same formulas give
$W=O(x^{2\nu_1+2})\to0$, since $\nu_1>-1$.
Notice that for $d>0$ and $s<0$, $W$ diverges to $+\infty$;
no vanishing boundary condition is imposed in that case.

For large positive arguments the classical expansions, including
their derivative expansions, give
\begin{equation}\label{bessel-eq:largeLog}
 \frac{d}{dx}\log I_\nu(ax)
 =a-\frac1{2x}+\frac{4\nu^2-1}{8ax^2}+O(x^{-3}).
\end{equation}
See \cite[\S10.40(i), Eqs.~10.40.1 and 10.40.3]{bessel-DLMF}.
If $a_1<a_2$, subtracting the two expansions gives the negative
limit $a_1-a_2$. If $a_1=a_2=a$ and $B\ne0$, it gives
\[
 \frac{f'}f-\frac{g'}g=\frac{B}{2ax^2}+O(x^{-3}).
\]
Since $W=xfg(f'/f-g'/g)$, positivity of $xfg$ proves the claims.
\end{proof}

\subsection{Proof of the classification}
\begin{proof}[Proof of Theorem~\ref{bessel-thm:main}]
The identical-parameter case is immediate. If $d=0$ and $D>0$,
then $B=0$ and \eqref{bessel-eq:identity} gives $W'<0$. Together with
$W(0+)=0$ from Lemma~\ref{bessel-lem:endpoints}, this yields $W<0$ on
$(0,\infty)$, hence strict decrease of $R$.

Now assume $d>0$. By Lemma~\ref{bessel-lem:endpoints}, $W$ is positive
near zero. First take $D>0$. If $B\le0$, then $W'<0$ everywhere.
If $B>0$, let $x_0=\sqrt{B/D}$. Then $W'>0$ on $(0,x_0)$ and
$W'<0$ on $(x_0,\infty)$. In particular $W(x_0)>0$.
In either case $W$ stays positive until it enters its strictly
decreasing region. It is negative for all sufficiently large $x$
by Lemma~\ref{bessel-lem:endpoints}. Continuity and strict decrease give
exactly one zero $x_*$, and its sign pattern is $(+,-)$.
Equation~\eqref{bessel-eq:identity} transfers that pattern to $R'$.
When $B>0$, the proof also shows $x_*>\sqrt{B/D}$.

Finally take $D=0$. If $s>0$, then $B=ds>0$ and $W'>0$,
so positivity near zero gives $W>0$ everywhere. If $s=0$, then
$W'=0$ and \eqref{bessel-eq:smallW} gives the strictly positive constant
$W=Cd$. These are case~4. If $s<0$, then $B<0$ and $W'<0$.
Now $W$ is positive near zero and eventually negative by
Lemma~\ref{bessel-lem:endpoints}. It has exactly one zero, proving case~5.
All cases in the statement have been covered.
\end{proof}

\subsection{Transfer to Nuttall \texorpdfstring{$Q$}{Q} ratios}
Define, whenever the integral converges,
\begin{equation}\label{bessel-eq:Q}
 Q_{\mu,\nu}(a,b)=\int_b^\infty
 t^\mu e^{-(t^2+a^2)/2}I_\nu(at)\,dt.
\end{equation}
For $b=0$ and $\nu>-1$, convergence holds exactly when
$\mu+\nu>-1$; for $b>0$ it holds for every real $\mu$.

\begin{theorem}\label{bessel-thm:Q}
Suppose $\nu_1\ge\nu_2>-1$, $0<a_1\le a_2$, and $b\ge0$.
Then
\[
 \mu\longmapsto
 \frac{Q_{\mu,\nu_1}(a_1,b)}{Q_{\mu,\nu_2}(a_2,b)}
\]
is unimodal on $(-1-\nu_2,\infty)$ if $b=0$, and on all of
$\R$ if $b>0$. If it is nonmonotone, its pattern is increasing
then decreasing. In particular the conclusion holds on $(0,\infty)$
without the restriction $\nu_1-\nu_2\in2\mathbb N$ in
\cite[Theorem~4.3]{bessel-KVZ}.
\end{theorem}
\begin{proof}
We use the established integral-ratio theorem
\cite[Theorem~2.6 and Example~1]{bessel-KVZ}: a strictly totally positive
kernel of order~3 preserves unimodality of the numerator/denominator
ratio under integration against a positive weight, provided both
integrals converge uniformly on compact subsets of the parameter
interval. The direction of a nonmonotone pattern is preserved when
all kernel minors have positive signs.

Apply this with kernel $K(\mu,t)=t^\mu$, positive weight
$w(t)=e^{-t^2/2}$, and
\[
 A(t)=e^{-a_1^2/2}I_{\nu_1}(a_1t),\qquad
 B(t)=e^{-a_2^2/2}I_{\nu_2}(a_2t).
\]
The Mellin kernel is strictly totally positive of every order, as
is the exponential kernel $e^{\mu\log t}$.
The ratio $A/B$ is a positive constant multiple of the ratio in
Theorem~\ref{bessel-thm:main}. Its restriction to $(b,\infty)$ is still
unimodal, with an increasing-then-decreasing pattern if it is
nonmonotone.

For completeness, the convergence condition holds uniformly on compact
parameter intervals. At infinity, \eqref{bessel-eq:largeLog} or the standard
Bessel asymptotic shows domination by a polynomial times
$e^{-t^2/2+a_2t}$, which is integrable. Near zero, for a compact
interval with lower endpoint $\mu_0>-1-\nu_2$, the integrands
are bounded by constants times $t^{\mu_0+\nu_j}$, which are
integrable for $j=1,2$. For $b>0$ there is no endpoint singularity.
Both denominators are strictly positive. The cited theorem therefore
applies on the full claimed domains.
\end{proof}

At $b=0$ one obtains the corresponding gamma--Kummer ratio directly,
using the identity
\[
 Q_{\mu,\nu}(a,0)=2^{(\mu-\nu-1)/2}a^\nu e^{-a^2/2}
 \frac{\Gamma((\mu+\nu+1)/2)}{\Gamma(\nu+1)}
 {}_1F_1\left(\frac{\mu+\nu+1}{2};\nu+1;\frac{a^2}{2}\right).
\]
This identity follows by integrating \eqref{bessel-eq:series} term by term;
all terms are nonnegative on the convergence domain.

\subsection{Why the additional log-concavity clause fails}
The positive-order log-concavity assertion must be kept separate from
unimodality. An explicit, formally certified failure is already
recorded in \cite{bessel-TT}. The following analytic observation explains
a whole range of such failures.

\begin{proposition}\label{bessel-prop:lc}
For every $0<\nu<1/2$ and $0<a_1<a_2$, the function
$I_\nu(a_1x)/I_\nu(a_2x)$ is strictly decreasing on $(0,\infty)$,
but its logarithm is strictly convex for all sufficiently large $x$.
In particular it is not log-concave on the entire positive half-line.
\end{proposition}
\begin{proof}
Strict decrease is case~2 of Theorem~\ref{bessel-thm:main}.
The standard differentiated Bessel expansions give, for fixed
$a_1,a_2,\nu$,
\[
 \frac{d^2}{dx^2}\log\frac{I_\nu(a_1x)}{I_\nu(a_2x)}
 =\frac{1-4\nu^2}{4x^3}
    \left(\frac1{a_1}-\frac1{a_2}\right)+O(x^{-4}).
\]
For instance this follows from \eqref{bessel-eq:largeLog} and the derivative
expansions in \cite[\S10.40(i)]{bessel-DLMF}; the Bessel differential equation
also supplies the second derivative expansion directly.
The leading coefficient is strictly positive under the stated
assumptions, giving the assertion.
\end{proof}

\subsection{Scope and methodological dependencies}
The complete result proved here is the unimodality clause of the
2025 conjecture, together with its Nuttall $Q$ transfer and the full
monotonicity classification in Theorem~\ref{bessel-thm:main}.
The exact global parameter region for log-concavity is not determined.
Neither numerical sampling nor a finite interval computation enters
any proof. The endpoint expansions are ordinary fixed-parameter
asymptotic theorems, used only to determine an eventual sign.

The effective combination is a second-order differential equation,
the Lagrange/weighted-Wronskian identity, endpoint asymptotics, and
total positivity of an integral transform. These are classical tools;
the manuscript does not claim that the identity or any foundational
method is newly invented. Independent mathematical
review and a fuller priority check remain appropriate before publication.

There is a concrete methodological bridge to the supplied research
corpus \cite{VMRC}. Its item~006, Section~6, eliminates two solutions
of a second-order recurrence to obtain the discrete Abel identity
$W_n=-(c_n/a_n)W_{n-1}$. Here the same elimination idea is applied
to two differential equations: the common first-derivative term
cancels and the difference of the two potentials gives
\eqref{bessel-eq:identity}. Item~146, in its curvature-to-slope-imbalance
lemma, transfers a controlled derivative sign to the sign of a
comparison quantity by fixing its boundary behavior. The argument
above similarly transfers the one-change sign pattern of $W'$ to
the zero pattern of $W$, then to the derivative of the desired ratio.
These are methodological correspondences; the original arithmetic
and free-stable theorems are not assumed in this proof.

\section{Exact densities for multiplicative four-point configurations}
\label{sec:density}
\subsection{Statement and conventions}
For $A\subseteq\N$, write $A(X)=|A\cap[1,X]|$ and
\[
\underline d(A)=\liminf_{X\to\infty}\frac{A(X)}X,
\qquad
\overline d_{\log}(A)=\limsup_{X\to\infty}
\frac1{\log X}\sum_{n\le X,\ n\in A}\frac1n.
\]
The natural density $d(A)$ is the limit of $A(X)/X$ when it exists.
We call $A$ \emph{$(a,b,c)$-free} if there is no positive integer $x$
for which all of $x,ax,bx,cx$ belong to $A$.
For an integer $q\ge2$, let
$\nu_q(n)=\max\{e\ge0:q^e\mid n\}$; primality of $q$ is not required.

\begin{theorem}\label{density-thm:main}
Let $a,b,c$ be pairwise coprime integers with $2\le a<b<c$. Put
\[
D(a,b,c)=
\begin{cases}
\displaystyle 1-\frac{ab+c}{(a+1)(b+1)(c+1)},&c\le a+b+1,\\[5pt]
\displaystyle \frac{c}{c+1},&c\ge a+b+1.
\end{cases}
\]
Every $(a,b,c)$-free set $A$ satisfies
\[
\sum_{n\le X,\ n\in A}\frac1n\le D(a,b,c)\log X+O_{a,b,c}(1).
\]
Consequently $\overline d_{\log}(A)\le D(a,b,c)$ and
$\underline d(A)\le D(a,b,c)$. The largest attainable natural
density equals $D(a,b,c)$.
For $c\le a+b+1$, an extremizer is
\[
 A_* =\{n\ge1:
 (\nu_a(n),\nu_b(n),\nu_c(n))\bmod2
 \notin\{(0,0,1),(1,1,0)\}\}.
\]
For $c\ge a+b+1$, one may instead take
$A_* =\{n:\nu_c(n)\text{ is even}\}$.
\end{theorem}

\begin{corollary}
The largest natural density of a set avoiding $\{x,2x,3x,5x\}$ is
\[
\boxed{\frac{61}{72}=0.8472222222\ldots.}
\]
The same number is the largest possible lower natural density and
upper logarithmic density. No claim about the largest upper natural
density is made here.
\end{corollary}

\subsection{The exact weighted corner problem}
For real $a,b,c>1$, put
\[
w(i,j,k)=a^{-i}b^{-j}c^{-k},\qquad
W=\sum_{i,j,k\ge0}w(i,j,k)
=\frac{a}{a-1}\frac{b}{b-1}\frac{c}{c-1}.
\]
For $e\in\Z_{\ge0}^3$, let
$\C_e=\{e,e+e_1,e+e_2,e+e_3\}$, where the $e_i$ are the standard
basis vectors. A transversal meets every $\C_e$.

\begin{lemma}[Two parity classes]\label{density-lem:code}
The set
\[
 T_* =\{(i,j,k):(i,j,k)\bmod2\in\{001,110\}\}
\]
meets each corner in exactly one vertex, and
\[
 \frac{w(T_*)}{W}
 =\frac{ab+c}{(a+1)(b+1)(c+1)}.
\]
\end{lemma}
\begin{proof}
Modulo two, each corner is a vertex of the binary three-cube together
with its three neighbors. The two antipodal vertices $001$ and $110$
have disjoint radius-one Hamming balls, each of size four; these balls
partition the cube. Hence every corner contains exactly one selected
parity class.
Under normalized geometric weights, the probability of an even
exponent in base $a$ is $a/(a+1)$, and that of an odd exponent is
$1/(a+1)$. Multiplication in the three coordinates gives the displayed
formula.
\end{proof}

\begin{theorem}[Exact weighted optimum]\label{density-thm:weighted}
Suppose $b\ge a+1$ and $b\le c\le a+b+1$. Then the minimum total weight of a
transversal of all corners is
\[
Y=W\frac{ab+c}{(a+1)(b+1)(c+1)},
\]
and $T_*$ attains it.
\end{theorem}
\begin{proof}
Define
\[
z_0=\frac1c,\qquad z_1=\frac{c-b}{(a+1)c},\qquad
 y(i,j,k)=
 \begin{cases}
 w(i,j,k)z_{j\bmod2},&k\text{ even},\\
 0,&k\text{ odd}.
 \end{cases}
\]
All these corner masses are nonnegative. The load of a vertex is
$\ell(e)=\sum_{f:e\in\C_f}y(f)$. If $k$ is odd, the only nonzero
contribution comes from $e-e_3$, and
\[
 \frac{\ell(i,j,k)}{w(i,j,k)}=c z_{j\bmod2}
 =\begin{cases}1,&j\text{ even},\\(c-b)/(a+1),&j\text{ odd}.
 \end{cases}
\]
These values are at most one. If $k$ is even, the exact formula is
\[
\frac{\ell(i,j,k)}{w(i,j,k)}
=(1+a\mathbf1_{i\ge1})z_{j\bmod2}
+b\mathbf1_{j\ge1}z_{1-j\bmod2}.
\]
For odd $j$, the largest possible load ratio is
$(a+1)z_1+bz_0=1$. For even $j$, it is at most
\[
(a+1)z_0+bz_1
=1+\frac{(b-a-1)(c-a-b-1)}{(a+1)c}\le1.
\]
The last inequality uses both hypotheses. Hence every vertex load
is at most its weight.
This calculation includes vertices on either coordinate boundary:
omitting an incoming corner only decreases the load.
Moreover the load equals the vertex weight at every vertex of $T_*$.
For class $001$, this follows from the odd-$k$, even-$j$ formula.
For class $110$, both $i$ and $j$ are positive, so the even-$k$
formula is exactly $(a+1)z_1+bz_0=1$.

For any transversal $T$, nonnegativity and Tonelli's theorem give
\[
 w(T)\ge\sum_{e\in T}\ell(e)
 =\sum_fy(f)|\C_f\cap T|\ge\sum_fy(f).
\]
For $T_*$ both inequalities are equalities, by the load calculation
and Lemma~\ref{density-lem:code}. Thus
$\sum_fy(f)=w(T_*)=Y$, proving the assertion.
\end{proof}

\begin{lemma}[Dominating last multiplier]\label{density-lem:dominant}
For real $a,b,c>1$ satisfying $c\ge1+a+b$, the minimum transversal
weight is $W/(c+1)$.
\end{lemma}
\begin{proof}
Use $y(i,j,k)=w(i,j,k)/c$ for even $k$ and zero for odd $k$.
At odd $k$ the load is $w$; at even $k$ the load is at most
$w(1+a+b)/c\le w$. The odd-$k$ vertices meet every corner, and
meet exactly once every corner assigned positive mass (those with
even base $k$). They have total weight $W/(c+1)$. The same primal--dual argument
proves optimality. This covers the entire second parameter region.
\end{proof}

\subsection{Transfer to arbitrary sets of positive integers}
We give the transfer in detail because a weighted lattice optimum
alone does not bound the upper natural density of arbitrary sets.
For pairwise coprime $a,b,c$, define
\[
M=\{m\ge1:a\nmid m,\ b\nmid m,\ c\nmid m\},\qquad
\delta=(1-a^{-1})(1-b^{-1})(1-c^{-1})=W^{-1}.
\]
Every positive integer has a unique representation
$n=ma^ib^jc^k$ with $m\in M$ and $i,j,k\ge0$.
The representation follows by repeatedly removing whole powers of
the three pairwise coprime bases. By inclusion--exclusion,
\begin{equation}\label{density-eq:Mh}
 \sum_{m\le X,\ m\in M}\frac1m=\delta\log X+O_{a,b,c}(1).
\end{equation}
For $z\ge1$, let
$\Lambda_z=\{(i,j,k):a^ib^jc^k\le z\}$.
Each corner is contained in $\Lambda_z$ exactly when its base
satisfies $a^ib^jc^k\le z/c$, since $c$ is the largest multiplier.
Let $y$ be the packing from Theorem~\ref{density-thm:weighted} or
Lemma~\ref{density-lem:dominant}, and write $Y=\sum_e y(e)$ and
\[
Y_z=\sum_{\C_e\subseteq\Lambda_z}y(e),\qquad R(z)=Y-Y_z.
\]

\begin{lemma}[Packing tail]\label{density-lem:tail}
For a constant depending only on the three bases,
\[
0\le R(z)\ll_{a,b,c}\frac{(1+\log z)^2}{z}\qquad(z\ge1).
\]
\end{lemma}
\begin{proof}
The masses $y(e)$ are at most a fixed constant times $w(e)$.
For $t\ge1$, the number of exponent triples with
$t<a^ib^jc^k\le2t$ is $O((1+\log t)^2)$: there are
$O((1+\log t)^2)$ possibilities for $(i,j)$, and for each pair
at most $1+\log_c2$ possible integers $k$.
Consequently the total weight above a threshold $u\ge1$ is bounded
by
\[
\sum_{h\ge0}\frac{O((1+\log(2^{h+1}u))^2)}{2^hu}
\ll\frac{(1+\log u)^2}{u}.
\]
Apply this with $u=z/c$ if $z\ge c$. The remaining bounded interval
$1\le z<c$ is absorbed by increasing the constant.
\end{proof}

\begin{proof}[Upper bound in Theorem~\ref{density-thm:main}]
For each $m\in M$, the exponent vectors of members of $A$ form a
corner-free subset. Its complement in $\Lambda_{X/m}$ is therefore
a transversal of all corners contained in that finite region.
Restricting the nonnegative packing to these corners preserves the
load bound. The weighted mass of $A$ in the fiber is consequently
at most $w(\Lambda_{X/m})-Y+R(X/m)$.
Summing over fibers yields
\[
\sum_{n\le X,\ n\in A}\frac1n
\le \sum_{n\le X}\frac1n
-Y\sum_{m\le X,\ m\in M}\frac1m
+\sum_{m\le X}\frac{R(X/m)}m.
\]
By Lemma~\ref{density-lem:tail}, the last term is at most a constant times
\[
\frac1X\sum_{m\le X}(1+\log(X/m))^2=O(1).
\]
For example, splitting $m$ into the intervals
$(X2^{-h-1},X2^{-h}]$ bounds this expression by a convergent
constant multiple of $\sum_{h\ge0}2^{-h}(h+1)^2$.
Equation~\eqref{density-eq:Mh} now gives
\[
\sum_{n\le X,\ n\in A}\frac1n
\le(1-\delta Y)\log X+O(1).
\]
The two weighted optimum formulas give precisely the two cases of
$D(a,b,c)$.
Finally partial summation shows
$\underline d(A)\le\overline d_{\log}(A)$: if
$A(t)\ge(d_0-\varepsilon)t$ for all sufficiently large $t$, then
$\sum_{n\le X,n\in A}1/n\ge(d_0-\varepsilon)\log X+O(1)$.
Taking $d_0=\underline d(A)$ proves the stated density consequences.
\end{proof}

\begin{proof}[Attainment in Theorem~\ref{density-thm:main}]
Multiplication by each base increments exactly its own exponent,
because the bases are pairwise coprime. Lemma~\ref{density-lem:code} thus
shows that $A_*$ is $(a,b,c)$-free.
For every fixed exponent triple, the integers with that triple
have density $\delta a^{-i}b^{-j}c^{-k}$, by inclusion--exclusion.
Summing these densities is justified uniformly: the union of triples
with any exponent at least $L$ has upper density at most
$a^{-L}+b^{-L}+c^{-L}$, which tends to zero.
Therefore the complement of $A_*$ has density
$\delta w(T_*)=(ab+c)/((a+1)(b+1)(c+1))$.
In the dominating case the even-$\nu_c$ construction is free since
$x$ and $cx$ cannot both belong to it, and its density is
$(1-c^{-1})\sum_{t\ge0}c^{-2t}=c/(c+1)$.
This also proves attainment for lower natural and upper logarithmic
densities, since a set with natural density has the same logarithmic
density.
\end{proof}

\subsection{The example and proof checks}
For $(a,b,c)=(2,3,5)$ the corner masses are
\[
y(i,j,k)=
\begin{cases}
2^{-i}3^{-j}5^{-k}/5,&j,k\text{ even},\\
2\cdot2^{-i}3^{-j}5^{-k}/15,&j\text{ odd},\ k\text{ even},\\
0,&k\text{ odd}.
\end{cases}
\]
The excluded parity classes have densities
\[
\Pr(001)=\frac23\frac34\frac16=\frac1{12},\qquad
\Pr(110)=\frac13\frac14\frac56=\frac5{72}.
\]
Their sum is $11/72$.
The naive even-$\nu_5$ construction has density $5/6$; the exact
optimum exceeds it by $1/72$. Other examples are
\[
D(3,4,5)=\frac{103}{120},\quad
D(3,4,7)=\frac{141}{160},\quad
D(2,5,7)=\frac{127}{144},\quad
D(3,5,7)=\frac{85}{96}.
\]
The standard-library verifier in Appendix~\ref{app:density-code} enumerates all $256$
binary parity patterns, retaining exactly those that hit all eight
parity corners. For every pairwise coprime integer triple
$2\le a<b<c\le20$, it computes the least weight of every retained
pattern with exact rational arithmetic and compares that independently
computed value with the theorem. It also checks all boundary/parity
load types. These checks supplement the displayed symbolic proof
and are not a finite substitute for it.

\subsection{Provenance and limits}
The research connection is the product-weighted lattice packing
method used in corpus item [140], combined with the elementary
binary perfect code $\{001,110\}$. The code itself is classical.
The general setting of extremal density and multiplicative linear
forms appears in Graham--Spencer--Witsenhausen (1977) and
Chung--Erd\H{o}s--Graham (2002). Veselinov (2026) considers the related
finite and smooth-number extremal problems. The present classification
is about the natural density of an infinite set; finite-prefix extrema
and upper natural density are different quantities. A bounded search
did not establish whether this exact formula had appeared previously.
We do not label it a resolution of a named conjecture, establish priority,
or claim external peer review. The proof is entirely elementary and
does not depend on any unverified assertion from the corpus.
The largest upper natural density is a different extremal problem and
is not calculated here. No irrationality conclusion is drawn.


\clearpage
\section{Synthesis, dependencies, and verification}
The common feature is a proof transfer with a controlled sign or weight:
\begin{itemize}
\item \textbf{Graph spectra.} The external square-energy inequality
supplies a lower bound on spectral mass. An eigenvector sign-support
bound makes the resulting inequality strict, and the classical
Favaron--Mah\'eo--Sacl\'e bound transfers it to the Randi\'c index.
\item \textbf{Special functions.} An exact differential identity controls
the weighted Wronskian. Endpoint signs determine its zero pattern;
the established integral-ratio theorem then transfers unimodality.
\item \textbf{Multiplicative density.} Nonnegative corner masses give a
dual packing bound, with equality on explicit parity classes. A tail
estimate transfers the lattice optimum to arbitrary integer sets.
\end{itemize}
The graph theorem does not assume the corpus radius theorem. The Bessel
argument does not assume the corpus Franel or free-stable conclusions.
The density proof does not assume an unverified corpus assertion.
The original corpus connections are methodological and remain identified
in the relevant sections and in the consolidated reference~\cite{VMRC}.

The bounds on density concern natural density, lower natural density,
and upper logarithmic density. Maximum upper natural density is not
determined. The Bessel classification does not supply a complete
log-concavity parameter region. Historical priority and external
acceptance are not established by this editorial consolidation.

The original finite-check programs are reproduced in the appendices.
The graph program uses floating-point eigenvalues and is a sanity check
only. The arithmetic program uses exact rational arithmetic on a finite
family. Neither is used as a premise of any theorem in this manuscript.
The two programs were rerun for this edition; their printed results are
recorded alongside the complete sources.

\appendix
\clearpage
\section{Graph sanity-check program}\label{app:graph-code}
This program requires Python, NumPy and NetworkX. It checks the connected
graphs of orders two through seven in the NetworkX graph atlas. The
analytic proofs in Section~\ref{sec:graph} apply to all finite orders.
\begin{lstlisting}
"""Floating-point sanity checks only; the manuscript contains the proof.

Checks every connected graph of orders 2 through 7 in NetworkX's graph atlas.
Dependencies: networkx and numpy. No numerical output is a proof certificate.
"""
import json
import math
import networkx as nx
import numpy as np

count = 0
quadratic_count = 0
equality_count = 0
minimum_gap = float('inf')
for graph in nx.graph_atlas_g():
    n = len(graph)
    if n < 2 or not nx.is_connected(graph):
        continue
    count += 1
    m = graph.number_of_edges()
    values = np.linalg.eigvalsh(nx.to_numpy_array(graph))
    largest, second = values[-1], values[-2]
    randic = sum(1 / math.sqrt(graph.degree(u) * graph.degree(v))
                 for u, v in graph.edges())
    assert randic > second + 1 + 1e-9
    assert 2*m > largest**2 + (second+1)**2 + 1e-9
    minimum_gap = min(minimum_gap, randic-second-1)
    if n >= 4 or m < n*(n-1)//2:
        quadratic_count += 1
        gap = randic**2-second**2-n+1
        assert gap > -1e-8
        star = sorted(dict(graph.degree()).values()) == [1]*(n-1)+[n-1]
        complete_four = n == 4 and m == 6
        expected_equality = star or complete_four
        assert (abs(gap) < 1e-8) == expected_equality
        equality_count += int(expected_equality)
print(json.dumps({'status':'PASS (floating-point sanity checks, not proof)',
                  'connected_graphs':count,
                  'quadratic_bound_graphs':quadratic_count,
                  'predicted_equality_graphs':equality_count,
                  'smallest_observed_randic_minus_lambda2_minus_one':minimum_gap},
                 indent=2))
\end{lstlisting}
\noindent\textbf{Recorded output.}
\begin{lstlisting}[language={}]
{
  "status": "PASS (floating-point sanity checks, not proof)",
  "connected_graphs": 995,
  "quadratic_bound_graphs": 993,
  "predicted_equality_graphs": 6,
  "smallest_observed_randic_minus_lambda2_minus_one": 0.23427568761990858
}
\end{lstlisting}

\clearpage
\section{Exact arithmetic verification program}\label{app:density-code}
This standard-library Python program checks parity covers and the
packing loads, including coordinate boundaries, for the coprime triples
in the displayed finite range. Section~\ref{sec:density} supplies the
symbolic proof for all admissible triples.
\begin{lstlisting}
#!/usr/bin/env python3
"""Exact independent finite checks for the all-triples density theorem.

No packages required. This is a sanity check supplementing the symbolic proof,
not a proof by checking finitely many integer triples.
"""
from fractions import Fraction as Q
from itertools import combinations, product
from math import gcd

vertices=list(product(range(2),repeat=3))
corners=[sum(1<<u for u in (v,v^4,v^2,v^1)) for v in range(8)]
feasible=[mask for mask in range(256) if all(mask&corner for corner in corners)]
code=(1<<1)|(1<<6)
assert code in feasible
assert all((code&corner).bit_count()==1 for corner in corners)

triples=loads=0
for a,b,c in combinations(range(2,21),3):
    if any(gcd(u,v)>1 for u,v in combinations((a,b,c),2)):
        continue
    triples+=1
    den=(a+1)*(b+1)*(c+1)
    weights=[Q(a**(1-i)*b**(1-j)*c**(1-k),den) for i,j,k in vertices]
    assert sum(weights)==1
    brute=min(sum(weights[v] for v in range(8) if mask>>v&1)
              for mask in feasible)
    code_mass=Q(a*b+c,den)
    parity_mass=Q(1,c+1)
    predicted=min(code_mass,parity_mass)
    assert brute==predicted,(a,b,c,brute,predicted)
    assert (code_mass<=parity_mass)==(c<=a+b+1)

    def w(i,j,k):return Q(1,a**i*b**j*c**k)
    def y(i,j,k):
        if min(i,j,k)<0 or k%2:return Q(0)
        if c<=a+b+1:
            z=Q(1,c) if j%2==0 else Q(c-b,(a+1)*c)
        else:z=Q(1,c)
        return w(i,j,k)*z
    for i,j,k in product(range(3),repeat=3):
        loads+=1
        load=y(i,j,k)+y(i-1,j,k)+y(i,j-1,k)+y(i,j,k-1)
        assert 0<=load<=w(i,j,k),(a,b,c,i,j,k,load,w(i,j,k))
        selected=((i%2,j%2,k%2) in ((0,0,1),(1,1,0))
                  if c<=a+b+1 else k%2==1)
        if selected:assert load==w(i,j,k)

examples={(2,3,5):Q(61,72),(3,4,5):Q(103,120),
          (3,4,7):Q(141,160),(2,5,7):Q(127,144),(3,5,7):Q(85,96)}
for (a,b,c),expected in examples.items():
    actual=1-min(Q(a*b+c,(a+1)*(b+1)*(c+1)),Q(1,c+1))
    assert actual==expected
    print('D',a,b,c,'=',actual)
print('PASS:',len(feasible),'admissible parity patterns;',triples,
      'coprime triples;',loads,'exact boundary/parity load checks')
\end{lstlisting}
\noindent\textbf{Recorded output.}
\begin{lstlisting}[language={}]
D 2 3 5 = 61/72
D 3 4 5 = 103/120
D 3 4 7 = 141/160
D 2 5 7 = 127/144
D 3 5 7 = 85/96
PASS: 183 admissible parity patterns; 297 coprime triples; 8019 exact boundary/parity load checks
\end{lstlisting}

\addcontentsline{toc}{section}{References}
\begin{thebibliography}{99}\small\raggedright
\bibitem{graph-Caz}
T.~Cazenave,
\emph{Lemmas and Conjectures in Spectral Graph Theory},
arXiv:2611.03879v1 (2026), Section~7, Conjecture~1 (\#643),
and Section~2.2, Theorem~2.
\url{https://arxiv.org/html/2611.03879v1}.

\bibitem{graph-LTZ}
Y.~Liu, Q.~Tang and S.~Zhang,
\emph{The positive and negative square-energy conjecture},
arXiv:2607.18031v1 (2026), Theorems~1.2 and~2.1.
\url{https://arxiv.org/html/2607.18031v1}.

\bibitem{graph-FMS}
O.~Favaron, M.~Mah\'eo and J.-F.~Sacl\'e,
\emph{Some eigenvalue properties in graphs (conjectures of Graffiti II)},
Discrete Mathematics \textbf{111} (1993), 197--220.

\bibitem{bessel-KVZ}
D. Karp, A. Vishnyakova and Y. Zhang,
\emph{Unimodality Preservation by Ratios of Functional Series and
Integral Transforms}, Results Math. \textbf{80}, 112 (2025).
\href{https://doi.org/10.1007/s00025-025-02431-4}{doi:10.1007/s00025-025-02431-4}.
Author preprint \href{https://arxiv.org/html/2408.01755v2}{arXiv:2408.01755v2};
Theorems~2.6 and~4.3, Example~1, Conjecture~1.

\bibitem{bessel-DLMF}
NIST Digital Library of Mathematical Functions,
\href{https://dlmf.nist.gov/10.25}{\S10.25} and
\href{https://dlmf.nist.gov/10.40}{\S10.40}.

\bibitem{bessel-Segura}
J. Segura, \emph{Monotonicity Properties for Ratios and Products of
Modified Bessel Functions and Sharp Trigonometric Bounds},
Results Math. (2021).
\href{https://doi.org/10.1007/s00025-021-01531-1}{doi:10.1007/s00025-021-01531-1}.
The adjacent-order Riccati framework is related background.

\bibitem{bessel-TT}
Trureturing, \emph{A modified Bessel ratio need not be log-concave},
formal proof record, consulted 15 November 2026.
\url{https://www.trureturing.org/research/karp-vishnyakova-zhang-2025-bessel-ratio-log-concavity/}.
This source concerns the negative log-concavity clause, not the
positive unimodality proof above.

\bibitem{density-GSW}
R. L. Graham, J. H. Spencer and H. S. Witsenhausen,
\emph{On extremal density theorems for linear forms}, in
Number Theory and Algebra, Academic Press, 1977, pp. 103--109.

\bibitem{density-CEG}
F. Chung, P. Erd\H{o}s and R. Graham,
\emph{On sparse sets hitting linear forms}, in
Number Theory for the Millennium, Vol. I, 2002, pp. 257--272.

\bibitem{density-Ves}
N. Veselinov, \emph{Extremal densities for forbidden
configurations in $S$-smooth numbers}, arXiv:2604.15515v1 (2026),
\url{https://arxiv.org/abs/2604.15515}.

\bibitem{VMRC}
\emph{Verifiable Mathematical Reasoning Corpus}, research drafts (2026).
Methodological references in this manuscript: item~006,
\emph{Minimal Recurrence Orders and Integral Solution Lattices for Franel Sequences};
item~140, \emph{Erd\H{o}s--Graham Density, Parity Packings};
item~146, \emph{Perfect Skewness of Free Stable Densities}; and
item~161, \emph{Randi\'c Index and Vertex Radius}.
The density source cites repository snapshot
\texttt{38348353416f1f7fc50cecb93c3b554a00260c13}.
\url{https://github.com/QEEDDD/Verifiable-Mathematical-Reasoning-Corpus}.
\end{thebibliography}
\end{document}
